% GENERATED FILE. Edit canonical lessons, then run npm run generate:books.
% canonical-source-hash: fnv1a64:d65caacb90a177fd
% canonical-lessons: LA-C451-percontor, LA-C451-sciscitor, LA-C451-requiro, LA-C451-monstro, LA-C451-indico

\chapter{To inquire, To ask around, To ask after, To show, To make known}
\label{ch:la-to-inquire-to-ask-around-to-ask-after-to-show-to-make-known}
\hlchaptermodality{451}

\begin{chapteropening}
\textbf{By the end of this chapter:} \emph{I can ask the way and be shown it: inquiring, asking around, asking after a place, showing and making known.}

Complete the last lesson of chapter 451: I can ask the way and be shown it: inquiring, asking around, asking after a place, showing and making known.
\end{chapteropening}

\section[percontor, percontārī]{percontor, percontārī — to inquire, to ask questions}

\label{lesson:LA-C451-percontor}



\begin{quote}
Before the new one: say the Latin for grasping, then the Latin for dirt.
\end{quote}

\subsection*{You'll want to know: percontor, percontārī}
\textbf{percontor, percontārī} — \textquotedblleft{}to inquire, to ask questions\textquotedblright{}.

\begin{quote}
Viātor ā pastōre percontātur ubi oppidum sit.
\end{quote}

\begin{quote}
\emph{The traveller asks a shepherd where the town is.}
\end{quote}

Deponent: \emph{percontor, percontārī, percontātus sum}, with \emph{ā} or \emph{ex} and the person asked. An indirect question follows in the subjunctive: \emph{ubi sit}.

\begin{cousinweb}
The origin of \emph{percontor} is not settled. Lewis \& Short build it on \emph{per} and a root they connect with hesitating; the older spelling is \emph{percunctor}. EtymDB lists it without a parent.
\end{cousinweb}

\begin{grammarlens}[title={What you've built}]
One more word for telling, explaining and describing at B1.
\end{grammarlens}

\subsection*{Your turn}
\begin{itemize}
\raggedright
  \item \emph{Say it:} \emph{percontor, percontārī}
  \item \emph{Say it:} \emph{percontor, percontārī}, once more
  \item \emph{Say it:} a sentence of your own with \emph{percontor, percontārī}
  \item \emph{Recall:} say the Latin for grasping, then the Latin for dirt, then say \emph{percontor, percontārī} again
\end{itemize}

\begin{tcolorbox}[breakable,colback=teal!4,colframe=teal!35!black,title={Before you move on}]
What does percontor, percontārī mean? (\textquotedblleft{}To inquire, to ask questions\textquotedblright{}.) Say it once more.
\end{tcolorbox}
\section[scīscitor, scīscitārī]{scīscitor, scīscitārī — to ask around, to try to find out}

\label{lesson:LA-C451-sciscitor}



\begin{quote}
Before the new one: say the Latin for dirt, then the Latin for to inquire.
\end{quote}

\subsection*{You'll want to know: scīscitor, scīscitārī}
\textbf{scīscitor, scīscitārī} — \textquotedblleft{}to ask around, to try to find out\textquotedblright{}.

\begin{quote}
Mercātor in forō dē viā scīscitātur.
\end{quote}

\begin{quote}
\emph{The merchant asks around in the forum about the road.}
\end{quote}

Deponent: \emph{scīscitor, scīscitārī, scīscitātus sum}. A frequentative: asking again and again.

\begin{cousinweb}
\emph{scīscitor} is built on \emph{scīscō} 'to seek to know' with the frequentative ending \emph{-itor} (EtymDB), and \emph{scīscō} comes from \emph{sciō} 'to know' (Lewis \& Short).
\end{cousinweb}

\begin{grammarlens}[title={What you've built}]
One more word for telling, explaining and describing at B1.
\end{grammarlens}

\subsection*{Your turn}
\begin{itemize}
\raggedright
  \item \emph{Say it:} \emph{scīscitor, scīscitārī}
  \item \emph{Say it:} \emph{scīscitor, scīscitārī}, once more
  \item \emph{Say it:} a sentence of your own with \emph{scīscitor, scīscitārī}
  \item \emph{Recall:} say the Latin for dirt, then the Latin for to inquire, then say \emph{scīscitor, scīscitārī} again
\end{itemize}

\begin{tcolorbox}[breakable,colback=teal!4,colframe=teal!35!black,title={Before you move on}]
What does scīscitor, scīscitārī mean? (\textquotedblleft{}To ask around, to try to find out\textquotedblright{}.) Say it once more.
\end{tcolorbox}
\section[requīrō, requīrere]{requīrō, requīrere — to ask after, to look for; to need}

\label{lesson:LA-C451-requiro}



\begin{quote}
Before the new one: say the Latin for to inquire, then the Latin for to ask around.
\end{quote}

\subsection*{You'll want to know: requīrō, requīrere}
\textbf{requīrō, requīrere} — \textquotedblleft{}to ask after, to look for; to need\textquotedblright{}.

\begin{quote}
Hospes caupōnam bonam requīrit.
\end{quote}

\begin{quote}
\emph{The visitor asks after a good inn.}
\end{quote}

\emph{requīrō, requīrere, requīsīvī, requīsītum}. English \emph{require} and \emph{request} come from it.

\begin{cousinweb}
\emph{requīrō} is built from \emph{re-} 'again' and \emph{quaerō} 'to seek, to ask' (Lewis \& Short), with the vowel weakened in the compound. English \emph{require} and \emph{request} come from it.
\end{cousinweb}

\begin{grammarlens}[title={What you've built}]
One more word for telling, explaining and describing at B1.
\end{grammarlens}

\subsection*{Your turn}
\begin{itemize}
\raggedright
  \item \emph{Say it:} \emph{requīrō, requīrere}
  \item \emph{Say it:} \emph{requīrō, requīrere}, once more
  \item \emph{Say it:} a sentence of your own with \emph{requīrō, requīrere}
  \item \emph{Recall:} say the Latin for to inquire, then the Latin for to ask around, then say \emph{requīrō, requīrere} again
\end{itemize}

\begin{tcolorbox}[breakable,colback=teal!4,colframe=teal!35!black,title={Before you move on}]
What does requīrō, requīrere mean? (\textquotedblleft{}To ask after, to look for; to need\textquotedblright{}.) Say it once more.
\end{tcolorbox}
\section[mōnstrō, mōnstrāre]{mōnstrō, mōnstrāre — to show, to point out}

\label{lesson:LA-C451-monstro}



\begin{quote}
Before the new one: say the Latin for to ask around, then the Latin for to ask after.
\end{quote}

\subsection*{You'll want to know: mōnstrō, mōnstrāre}
\textbf{mōnstrō, mōnstrāre} — \textquotedblleft{}to show, to point out\textquotedblright{}.

\begin{quote}
Agricola viātōribus viam ad flūmen mōnstrat.
\end{quote}

\begin{quote}
\emph{The farmer shows the travellers the way to the river.}
\end{quote}

\emph{mōnstrō, mōnstrāre, mōnstrāvī, mōnstrātum}, with a dative of the person shown. English \emph{demonstrate} comes from the compound, and \emph{muster} comes through French.

\begin{cousinweb}
\emph{mōnstrō} is built on \emph{mōnstrum} 'a sign, a portent' (EtymDB), which Lewis \& Short derive from \emph{moneō} 'to warn, to remind'. English \emph{demonstrate} comes from \emph{dēmōnstrāre}, and \emph{muster} comes from \emph{mōnstrāre} through French.
\end{cousinweb}

\begin{grammarlens}[title={What you've built}]
One more word for telling, explaining and describing at B1.
\end{grammarlens}

\subsection*{Your turn}
\begin{itemize}
\raggedright
  \item \emph{Say it:} \emph{mōnstrō, mōnstrāre}
  \item \emph{Say it:} \emph{mōnstrō, mōnstrāre}, once more
  \item \emph{Say it:} a sentence of your own with \emph{mōnstrō, mōnstrāre}
  \item \emph{Recall:} say the Latin for to ask around, then the Latin for to ask after, then say \emph{mōnstrō, mōnstrāre} again
\end{itemize}

\begin{tcolorbox}[breakable,colback=teal!4,colframe=teal!35!black,title={Before you move on}]
What does mōnstrō, mōnstrāre mean? (\textquotedblleft{}To show, to point out\textquotedblright{}.) Say it once more.
\end{tcolorbox}
\section[indicō, indicāre]{indicō, indicāre — to make known, to indicate, to show}

\label{lesson:LA-C451-indico}



\begin{quote}
Before the new one: say the Latin for to ask after, then the Latin for to show.
\end{quote}

\subsection*{You'll want to know: indicō, indicāre}
\textbf{indicō, indicāre} — \textquotedblleft{}to make known, to indicate, to show\textquotedblright{}.

\begin{quote}
Mīliārium indicat urbem decem mīlia passuum abesse.
\end{quote}

\begin{quote}
\emph{The milestone shows that the city is ten miles away.}
\end{quote}

\emph{indicō, indicāre, indicāvī, indicātum}, first conjugation. Not the same as \emph{indīcō, indīcere} 'to proclaim'. English \emph{indicate} comes from it.

\begin{cousinweb}
\emph{indicō} is built from \emph{in-} and the root of \emph{dīcō} 'to say, to show' (EtymDB; Lewis \& Short); \emph{index} 'a pointer' belongs to the same family. English \emph{indicate} comes from it.
\end{cousinweb}

\begin{grammarlens}[title={What you've built}]
One more word for telling, explaining and describing at B1.
\end{grammarlens}

\subsection*{Your turn}
\begin{itemize}
\raggedright
  \item \emph{Say it:} \emph{indicō, indicāre}
  \item \emph{Say it:} \emph{indicō, indicāre}, once more
  \item \emph{Say it:} a sentence of your own with \emph{indicō, indicāre}
  \item \emph{Recall:} say the Latin for to ask after, then the Latin for to show, then say \emph{indicō, indicāre} again
\end{itemize}

\begin{tcolorbox}[breakable,colback=teal!4,colframe=teal!35!black,title={Before you move on}]
What does indicō, indicāre mean? (\textquotedblleft{}To make known, to indicate, to show\textquotedblright{}.) Say it once more.
\end{tcolorbox}
